\section{Kurzfassung}
Ärzt*innen verwenden oft komplexe Begriffe bei der Diagnostik, wobei Laien meist nicht in der Lage sind, diesen Fachjargon zu verstehen und zu visualisieren.
Zudem sind Überweisungen zu Folgtetermine häufig notwendig, wodurch für die medizinische Einrichtung und die Patienten eine zusätzliche Belastung entsteht. 
Das Ziel dieser Diplomarbeit ist die Entwicklung eines Systems, das Patient:innen ermöglicht, sich mit dem:der Ärzt:in digitale Treffen auszumachen, um zuvor aufgenommene
MRT-Bilder gemeinsam zu betrachten und zu besprechen mittels VR-Brillen.